\chapter{Essential Tremor}
\index{Essential Tremor}
\label{sec:Essential Tremor}

\section{Overview}
Essential tremor is one of the most common movement disorders, affecting 1\% of the general population and 5\% of individuals over 65 years. This neurological disorder affects the central or peripheral nervous system, leading to progressive or episodic symptoms. It significantly impacts quality of life and often requires multidisciplinary care. Neurological disorders affect the central or peripheral nervous system, leading to a wide range of symptoms affecting movement, sensation, cognition, and autonomic function. The nervous system's complexity means that even small disruptions can cause significant functional impairment. Many neurological conditions are chronic and progressive, requiring long-term management and rehabilitation. Advances in neuroimaging and molecular diagnostics have improved understanding and treatment of these conditions.
\section{Causes}
This condition happens because of abnormal oscillatory activity in the cerebello-thalamo-cortical circuit, with genetic heterogeneity (autosomal dominant inheritance in familial cases, ETM1/ETM2 loci).    Protein aggregation and accumulation is a common mechanism in many neurodegenerative disorders.  Genetic mutations account for some neurological conditions, while others are sporadic or environmentally triggered.   
\section{Risk Factors}
\begin{itemize}[noitemsep]
  \item Genetic predisposition and family history
  \item Advancing age
  \item Lifestyle factors (diet, exercise, smoking, alcohol)
  \item Environmental exposures
  \item Underlying medical conditions
  \item Medication use
\end{itemize}
\section{Signs and Symptoms}
\subsection{Common symptoms}
\begin{itemize}[noitemsep]
  \item You may experience bilateral, symmetric postural and kinetic tremor of the upper limbs (6--12 Hz), which may also involve the head (titubation), voice, chin, or trunk. Tremor is typically absent at rest, worsens with goal-directed movement, and is exacerbated by stress, fatigue, and caffeine. Alcohol consumption temporarily reduces tremor amplitude in 50--70\% of patients. Headache and facial pain Weakness or paralysis of muscles Numbness, tingling, or abnormal sensations Vision changes including blurred or double vision Difficulty with balance and coordination Cognitive changes (memory loss, confusion, difficulty concentrating) Speech and language difficulties Seizures or involuntary movements Dizziness and vertigo Fatigue and sleep disturbances
  \item Pain or discomfort in the affected area
  \end{itemize}
\subsection{Emergency warning signs}
\begin{itemize}[noitemsep]
  \item Severe or worsening symptoms of essential tremor require immediate medical evaluation
\end{itemize}
\section{Progression and Stages}
The condition may progress through different stages of severity over time. Early stages often involve mild or subtle symptoms that may be overlooked. Moderate stages involve more noticeable symptoms affecting daily function. Advanced stages may involve significant impairment and complications.
\section{Complications}
\begin{itemize}[noitemsep]
  \item Disease progression leading to worsening symptoms
  \item Secondary infections or injuries
  \item Side effects of treatment
  \item Reduced quality of life
  \item Emotional and psychological impact
\end{itemize}
\section{Diagnosis}
Doctors usually diagnose this based on your symptoms and a physical exam, based on the Movement Disorder Society consensus criteria and exclusion of other causes (Parkinson disease, dystonic tremor, enhanced physiologic tremor, drug-induced tremor). Neurological diagnosis combines clinical history and examination findings with neuroimaging (MRI, CT), electrophysiological studies (EEG, EMG, nerve conduction), and laboratory testing of blood and cerebrospinal fluid.

\begin{itemize}[noitemsep]
  \item Comprehensive medical history and physical examination
  \item Laboratory tests to assess organ function
  \item Imaging studies to visualize affected structures
  \item Specialized testing depending on the affected system
  \item Consultation with relevant specialists
\end{itemize}
\section{Prevention}
\begin{itemize}[noitemsep]
  \item Maintain a healthy lifestyle with balanced diet and exercise
  \item Avoid known risk factors
  \item Attend regular medical check-ups for early detection
  \item Follow recommended screening guidelines
\end{itemize}
\section{Treatment Options}
Treatment options include The first-choice treatment with propranolol or primidone, alternative agents (topiramate, gabapentin, benzodiazepines), and for medication-refractory cases, deep brain stimulation of the ventral intermediate nucleus (VIM) of the thalamus or focused ultrasound thalamotomy. Neurological treatment focuses on symptom management, disease modification when possible, and rehabilitation to maintain function. A multidisciplinary team approach is often beneficial. Pharmacologic therapy targeting specific neurotransmitter systems or disease mechanisms Physical, occupational, and speech therapy for functional maintenance Surgical interventions when appropriate (deep brain stimulation, tumor resection) Cognitive rehabilitation and memory support Pain management for neuropathic pain Assistive devices and mobility aids Lifestyle modifications including exercise, nutrition, and sleep hygiene Support services and caregiver education Regular neurologic monitoring and treatment adjustment Clinical trials for emerging therapies

\begin{itemize}[noitemsep]
  \item Medications to manage symptoms and address underlying mechanisms
  \item Lifestyle modifications and self-care measures
  \item Physical therapy and rehabilitation
  \item Surgical intervention when indicated
  \item Regular monitoring and follow-up care
\end{itemize}
\section{Living with It}
\begin{itemize}[noitemsep]
  \item Follow your treatment plan as prescribed
  \item Maintain regular follow-up with your healthcare provider
  \item Make necessary lifestyle adjustments
  \item Seek support from family, friends, or support groups
  \item Stay informed about your condition
\end{itemize}
\section{Outlook}
\begin{itemize}[noitemsep]
  \item Outlook is non-cancerous but slowly progressive
  \item Disability from tremor may be significant.
\end{itemize}
